\section{Delphi：机器究竟在预测什么伦理对象}

第三项研究把 commonsense 扩展到 morally salient situations。这里风险显著上升：把多数人的描述性判断误写成“正确道德”会把 dataset bias 包装成权威。因此必须先定义 prediction target，再谈 accuracy。

\subsection{Machine Ethics 不是未来问题}

Delphi paper title slide 把系统定义为 experiment。紧接着的 slide 说神经语言模型大规模部署要求 machine ethics，并引用“机器人第一定律只是稻草人”式警告：真实系统不会遇到一个统一、可枚举的道德规则表，而是在内容过滤、排序、建议与交互中不断做隐含选择。

\lecturefigure{slide-24-delphi-paper.jpg}{Can Machines Learn Morality? The Delphi Experiment}{官方录像恢复幻灯片 V024；Jiang et al., 2021}

\lecturefigure{slide-25-machine-ethics-deployment.jpg}{广泛部署的 neural systems 使 machine ethics 成为现实约束}{官方录像恢复幻灯片 V025；Stanford CS25 Lecture 16}

\teachervoice{讲者不是因为 Delphi 完美才研究 machine ethics。她的论证恰好相反：当前系统已经在做带伦理后果的决定，如果研究者以“道德太难”为由退出，默认规则和既有偏差仍会继续运行。}

\subsection{Ask Delphi 的三种 Query Interface}

公开 demo slide 显示 free-form 输入、系统回答与显著 disclaimer。界面提醒“research prototype, still work in progress”，这是系统定义的一部分，不是法律页脚。它告诉使用者输出来自模型预测，不能作为真实行动、惩罚或裁决的唯一依据。

\lecturefigure{slide-26-ask-delphi-demo-disclaimer.jpg}{Ask Delphi demo 与研究原型免责声明}{官方录像恢复幻灯片 V026；Stanford CS25 Lecture 16}

下一页并排 free-form QA、relative QA 与 yes/no QA。free-form 评价单个情境，relative 比较两个行为，yes/no 回答特定判断。任务形式改变 label space，也改变攻击面：relative mode 很容易被构造为荒谬两害比较，课堂说明该功能在公开后数小时内就被下线。

\lecturefigure{slide-27-delphi-freeform-relative-yesno.jpg}{Delphi 的 free-form、relative 与 yes/no 判断接口}{官方录像恢复幻灯片 V027；Jiang et al., 2021}

\begin{importantbox}{Descriptive Prediction，不是 Normative Proof}
Delphi 学习的是“在训练数据覆盖的人群与情境中，人们通常如何评价”，可写成 $p(y\mid s,q)$。这不是从伦理公理推出 $y$，也不保证 minority view、公正原则或新文化情境被正确表达。
\end{importantbox}

\subsection{Compositional Situations 为什么更难}

组合 slide 从“在草坪上割草”逐步加入时间、地点、关系与例外条件。单个动作常有默认评价，但 context modifier 会翻转 judgment。读图时应比较同一 action 在不同 clause 下的变化；系统若只记住关键词，会忽略“半夜”“朋友工作时间”“紧急情况”等决定性条件。

\lecturefigure{slide-28-delphi-compositional-situations.jpg}{Delphi 对组合情境与上下文修饰的鲁棒性测试}{官方录像恢复幻灯片 V028；Jiang et al., 2021}

接着的 pair 展示 conventional scientific knowledge 与 common life know-how。模型知道“漂白剂与氨混合很危险”，却可能错误判断 lactose intolerance 的生活建议。这个对比说明百科事实、医学事实和 everyday norm 来自不同数据分布；一个总分会掩盖能力类型之间的差异。

\lecturefigure{slide-29-conventional-vs-common-life-knowledge.jpg}{Conventional scientific knowledge 与 common life know-how 的差别}{官方录像恢复幻灯片 V029；Stanford CS25 Lecture 16}

\subsection{Evaluation：先看 Task Family 再看总分}

结果图按 free-form、yes/no 与 relative QA 比较 GPT-3 zero-shot、few-shot、30-shot 与 Delphi。Delphi 在课堂展示的任务上更高，但这些结果依赖专门训练数据与 label definition。读图时不能把 supervised specialist 与 generic few-shot model 当成 architecture-only comparison；更合理的结论是 targeted moral data 明显改变了 judgment behavior。

\lecturefigure{slide-30-delphi-results-vs-gpt3.jpg}{Delphi 与 GPT-3 prompting baselines 的任务结果}{官方录像恢复幻灯片 V030；Jiang et al., 2021}

若 situation 为 $s$、query format 为 $q$、judgment label 为 $y$，训练目标可写为
\begin{equation}
\mathcal{L}_{\mathrm{Delphi}}(\theta)
=-\mathbb{E}_{(s,q,y)\sim\mathcal{D}_{\mathrm{norm}}}
\log p_{\theta}(y\mid s,q).
\end{equation}
$\mathcal{D}_{\mathrm{norm}}$ 是规范判断数据集合；这个公式直接说明模型学到什么取决于谁被采样、如何提问和怎样合并标签。

\subsection{本章小结}

Delphi 的 task definition 是描述性伦理判断。它在专门数据上比 generic prompting 更稳定，但 evaluation 不能把 accuracy 升格为 moral correctness。下一章必须打开数据与 backbone，理解这些 judgment 来自哪里。

